%Uncertainty quantification for natural language generation is a fast-growing area of research (see~\citet{baan2023uncertainty} and references therein). 
%With the rapid proliferation of LLMs and their swift adoption across various domains, uncertainty quantification for LLMs has garnered increasing attention. 
With the rapid proliferation of LLMs and their swift adoption across various domains, uncertainty quantification (UQ) for natural language generation (NLG) is a fast-growing area of research (see~\citet{baan2023uncertainty} and references therein). 
Adopting the language used by \citet{lin2023generating}, uncertainty measures the predictive distribution, while confidence (our focus) further depends on the specific generation.
A common approach involves reducing NLG to a de facto classification problem and utilizing or enhancing classical UQ methods for classification \citep{desai-durrett-2020-calibration,jiang-etal-2021-know,kamath-etal-2020-selective,wang-etal-2022-uncertainty,xiong2023llms}.
Recognizing the unique challenges posed by the inherently high (potentially infinite) dimensionality of NLG, recent research increasingly considers UQ for NLG from a sequence perspective~\citep{hou2023decomposing,kuhn2023semantic,malinin2021uncertainty,lin2023generating}.

Relatively less attention has been devoted to confidence measures. Sequence likelihood, or the log-probability of the generated sequence, remains one of the most popular proxies for assessing the quality of individual answers \citep{quach2023conformal,kuhn2023semantic,cole2023selectively}. Another natural approach, given the versatility and strong performance of LLMs, involves using prompts to elicit the LLM's own confidence level \cite{kadavath2022language,Lin2022TeachingMT,DBLP:journals/corr/abs-2012-14983,chen2023quantifying,he2023investigating,li2024think,wightman2023strength,tian-etal-2023-just}.
For free-form generation datasets, this is typically done by arranging answers as options and extracting the model's logits for each option. % \citep{he2023investigating}. 
While our method also employs prompts, they primarily guide the model to focus on tokens relevant to the context of the question. %, which we consider easier than directly asking the model for confidence. 
Another approach involves sampling additional generations and comparing their similarities~\citep{lin2023generating,cole2023selectively}, which demonstrates good discriminative capability but can be rather expensive. Although this approach has the advantage that it can work for black-box LLMs.
Ensemble methods have also been proposed~\citep{chen2023quantifying}.

Recently, \citet{duan2023shifting} proposed an improvement to sequence likelihood by weighting the tokens using their importance, which is computed by removing one token at a time from the original sequence and computing the NLI dissimilarity between the new and original sequences.
The removal of sub-word tokens, however, introduces drastically grammatically incorrect sentences, which could confuse the NLI model.
%Recently, \cite{duan2023shifting} proposed an approach to weight the logits of tokens differently. 
%It removes one token at a time from the original sequence, and use the dis-similarity between the new and original sequences as the importance weight for this token.
%This, however, introduces artificial non-words by removing sub-word tokens, resulting in sentences that can be drastically grammatically incorrect.
%Such out-of-distribution inputs could confuse the NLI model.
In addition, such a method could incur high computation overhead by making $n$ NLI comparisons, where $n$ is the length of the generation. 
In contrast, our method incurs no artificial grammatical errors and miniscule overhead. 
%In contrast, our method requires only one additional call to the base LLM. 
%Further, the soft attention it relies on is derived from the original generation, avoiding artificial grammatical errors.

An important downstream application of confidence measures is \textbf{abstention in LLMs}.
Classification with rejection~\cite{Corbiere2019AddressingConfidence,Fumera2000RejectThresholds,Geifman2017SelectiveNetworks,Jiang2018ToClassifier,lin2022scrib} could be considered the direct antecedent of selective NLG---both aiming to determine when to trust a model. % focus on finding conf score
Naturally, approaches similar to those in the classification with rejection literature have been applied to NLP applications~\cite{varshney-etal-2022-investigating,varshney-etal-2022-towards}.
More recently, \citet{cole2023selectively,yadkori2024mitigating,lin2023generating,quach2023conformal} started investigating NLG with abstention from UQ or risk control perspectives.

%\cite{yadkori2024mitigating}

Calibration is another crucial and relevant topic that has been extensively studied \cite{mielke-etal-2022-reducing,si-etal-2022-examining,xiong2023llms,zhu2023on}.
Although calibration is not directly related to our primary focus---our main interest lies in ranking the confidence of different measures---we include results on the calibrated performance in the Appendix.
Conformal prediction~\cite{vovk2005algorithmic} could also be considered a form of calibration at the distribution level, and has been extended to NLG to bound variants of error rates~\cite{quach2023conformal,yadkori2024mitigating}.
%\citet{yadkori2024mitigating} uses conformal prediction to bound the error rate (or ``hallucination rate'') of LLMs.

% From https://arxiv.org/pdf/2310.11732: Kadavath et al. (2022) show that advanced large pre-trained LMs are wellcalibrated while aligned LMs are mis-calibrated due to overconfidence on MCQs with logit-based UQ. Liang et al. (2023) suggests that the performance logit-based calibration of LMs depends on specific tasks and is orthogonal to other metrics such as accuracy and robustness.

% deng2024gotcha: needs fine-tuning
% he2023investigating: multiplace choice evlauation (https://arxiv.org/pdf/2310.11732)
% chen2024inside: eigenscore

% MC eval only: https://arxiv.org/pdf/2401.08694
